% problem_id: S046-epis-surj
% source_id: S046 (HoTT Book)
% source_file: setmath.tex
% statement_locator: setmath.tex lines 172--195
% proof_locator: setmath.tex lines 197--257
% env: lem
% pairing: tight (verified)
% status: complete (statement + author proof verbatim + reason + locators)
%
\begin{lem}\label{epis-surj}
For any function $f:A\to B$ between sets, the following are equivalent:
\begin{enumerate}
\item $f$ is an epimorphism.
\item Consider the pushout diagram
\begin{equation*}
  \xymatrix{
    {A}
    \ar[r]^{f}
    \ar[d]
    &
    {B}
    \ar[d]^{\iota}
    \\
    {\unit}
    \ar[r]_{t}
    &
    {C_f}
  }
\end{equation*}
in $\uset$ defining the mapping cone\index{cone!of a function}. Then the type $C_f$ is contractible.
\item $f$ is surjective.
\end{enumerate}
\end{lem}

\begin{proof}
Let $f:A\to B$ be a function between sets, and suppose it to be an epimorphism; we show $C_f$ is contractible.
The constructor $\unit\to C_f$ of $C_f$ gives us an element $t:C_f$.
We have to show that
\begin{equation*}
\prd{x:C_f} x= t.
\end{equation*}
Note that $x= t$ is a mere proposition, hence we can use induction on $C_f$.
Of course when $x$ is $t$ we have $\refl{t}:t=t$, so it suffices to find
\begin{align*}
I_0 & : \prd{b:B} \iota(b)= t\\
I_1 & : \prd{a:A} \opp{\alpha_1(a)} \ct I_0(f(a))=\refl{t}
\end{align*}
where $\iota:B\to C_f$ and $\alpha_1:\prd{a:A} \iota(f(a))= t$ are the other constructors
of $C_f$. Note that $\alpha_1$ is a homotopy from $\iota\circ f$ to
$\mathsf{const}_t\circ f$, so we find the elements
\begin{equation*}
\pairr{\iota,\refl{\iota\circ f}},\pairr{\mathsf{const}_t,\alpha_1}:
\sm{h:B\to C_f} \iota\circ f \htpy h\circ f.
\end{equation*}
By the dual of \cref{thm:mono}\ref{item:mono3} (and function extensionality), there is a path
\begin{equation*}
\gamma:\pairr{\iota,\refl{\iota\circ f}}=\pairr{\mathsf{const}_t,\alpha_1}.
\end{equation*}
Hence, we may define $I_0(b)\defeq \happly(\projpath1(\gamma),b):\iota(b)=t$.
We also have
\[\projpath2(\gamma) : \trans{\projpath1(\gamma)}{\refl{\iota\circ f}} = \alpha_1. \]
This transport involves precomposition with $f$, which commutes with $\happly$.
Thus, from transport in path types we obtain $I_0(f(a)) = \alpha_1(a)$ for any $a:A$, which gives us $I_1$.

Now suppose $C_f$ is contractible; we show $f$ is surjective.
We first construct a type family $P:C_f\to\prop$ by recursion on $C_f$, which is valid since \prop is a set.
On the point constructors, we define
\begin{align*}
P(t) & \defeq \unit\\
P(\iota(b)) & \defeq \brck{\hfiber{f}b}.
\end{align*}
To complete the construction of $P$, it remains to give a path
\narrowequation{\brck{\hfiber{f}{f(a)}} =_\prop \unit}
for all $a:A$.
However, $\brck{\hfiber{f}{f(a)}}$ is inhabited by $(f(a),\refl{f(a)})$.
Since it is a mere proposition, this means it is contractible --- and thus equivalent, hence equal, to \unit.
This completes the definition of $P$.
Now, since $C_f$ is assumed to be contractible, it follows that $P(x)$ is equivalent to $P(t)$ for any $x:C_f$.
In particular, $P(\iota(b))\jdeq \brck{\hfiber{f}b}$ is equivalent to $P(t)\jdeq \unit$ for each $b:B$, and hence contractible.
Thus, $f$ is surjective.

Finally, suppose $f:A\to B$ to be surjective, and consider a set $C$ and two functions
$g,h:B\to C$ with the property that $g\circ f = h\circ f$. Since $f$
is assumed to be surjective, for all $b:B$ the type $\brck{\hfib f b}$ is contractible.
Thus we have the following equivalences:
\begin{align*}
\prd{b:B} (g(b)= h(b))
& \eqvsym \prd{b:B} \Parens{\brck{\hfib f b} \to (g(b)= h(b))}\\
& \eqvsym \prd{b:B} \Parens{\hfib f b \to (g(b)= h(b))}\\
& \eqvsym \prd{b:B}{a:A}{p:f(a)= b} g(b)= h(b)\\
& \eqvsym \prd{a:A} g(f(a))= h(f(a))
\end{align*}
using on the second line the fact that $g(b)=h(b)$ is a mere proposition, since $C$ is a set.
But by assumption, there is an element of the latter type.
\end{proof}
